\originalchapter{22.02 作业四:势阱与测量}
题源: Spring 2012 Problem Set 4. 原题 1 标记课堂已解, 本次没有对应课堂解答文件; 本章解答均为编者补写. 理想测量与无限势阱之间的定义域问题明确列出.
\begin{exercise}\label{ex:2202-4-1}
原题 1. 势在 $a<x<a+L$ 为零, 区间外无限. (a) 测得位置 $a+L/13$ 后是什么态, $10\,\mu$s 后仍相同吗? (b) 再测动量, 有哪些结果, 之后是什么态、会演化吗? (c) 该态的位置均值? (d) 能量测得 $E'=3\hbar^2\pi^2/(2mL^2)$ 和 $E''=9\hbar^2\pi^2/(2mL^2)$ 的概率及平均能量?
\end{exercise}
\begin{solution}
(a) 理想形式态为 $\delta(x-a-L/13)$, 非归一化的 Hilbert 态, 具有无限动能尺度. 有限分辨率的窄波包才是物理态, 在无限阱内按离散能量基叠加演化, 通常 $10\,\mu$s 后不同; 要判断具体展宽需 m、L 与初宽.

(b) 保持硬墙不变时没有同 Dirichlet Hamiltonian 共同适用的自伴 ``平面波动量'' 投影测量, 原题未给测量模型. 若撤去墙并做自由动量测量, 所有实 p 都可出现, 形式位置本征态的 Fourier 模方常数且不可归一. 正规窄波包则有宽而可归一的动量谱. 测得 p 后自由广义态 $\e^{\ii px/\hbar}$ 只获相位.

(c),(d) 若课堂意图是把测后态形式上截断为阱内 $\psi_p=\e^{\ii p(x-a)/\hbar}/\sqrt L$, 可定义一个 L$^2$ 态来回答能量问题. 令 $\kappa=pL/\hbar$, 则位置均值 $a+L/2$, 能量展开系数
\begin{equation}
 c_n=\frac{\sqrt2n\pi[1-(-1)^n\e^{\ii\kappa}]}{(n\pi)^2-\kappa^2},\quad P(E_n)=|c_n|^2.
\end{equation}
$E'=3E_1$ 不在谱上, 概率零; $E''=9E_1=E_3$, 概率为 $18\pi^2|1+\e^{\ii\kappa}|^2/(9\pi^2-\kappa^2)^2$, 可去奇点按积分极限. 截断平面波在墙边非零, 不属于 H 的有限能量定义域, $\sum_nE_n|c_n|^2$ 发散. 因而不能同时声称它保持无限墙、动量确定且平均能量有限. 有限分辨率和实际测量耦合能消除这种理想化冲突; 原题缺此信息时以上条件回答比虚构一个有限值更明确.
\end{solution}
\begin{exercise}\label{ex:2202-4-2}
原题 2. (a) 什么是可观测量, 用什么对象表示? (b) H 有离散本征函数 $\psi_n$, 系统波函数是什么? (c) $\int w_4^*(x)\phi(x)\dd x$ 是测得 $E_4$ 的概率吗?
\end{exercise}
\begin{solution}
(a) 可测物理量用自伴算符及其谱测度表示. (b) 一般态为 $\sum_nc_n\psi_n$, $\sum|c_n|^2=1$; Hamiltonian 只给可能能级, 未给制备不能确定系数. (c) 该积分是概率振幅, 概率为模平方; 若简并需投影整个 $E_4$ 子空间.
\end{solution}
\begin{exercise}\label{ex:2202-4-3}
原题 3. 一维无限阱 $-L/2<x<L/2$. (a) 求全部归一能量态和能谱. (b) $\psi=a_2w_2+a_4w_4$ 时能量结果? (c) 给 $P_2=3/7,P_4=4/7$, 求系数. (d) 求在 $x=L/3$ 找到粒子的概率及位置均值. (e) 求动量结果与概率.
\end{exercise}
\begin{solution}
(a) $w_n=\sqrt{2/L}\sin[n\pi(x+L/2)/L]$, $E_n=n^2\hbar^2\pi^2/(2mL^2)$, $n\ge1$. (b) 只 $E_2,E_4$. (c) $a_2=\sqrt{3/7}\e^{\ii\alpha},a_4=\sqrt{4/7}\e^{\ii\beta}$, 能量概率未指定相对相位. (d) 精确单点概率零, 小区间 $\dd x$ 概率为 $|\psi(L/3)|^2\dd x$. 因 $w_2(L/3)=w_4(L/3)=-\sqrt{3/(2L)}$, 密度为 $3|a_2+a_4|^2/(2L)$. 两偶数 n 态均为奇宇称, 叠加密度为偶函数, 所以 $\langle x\rangle=0$.

(e) 若以自由动量表象测量 (如释放后飞行测量), 结果连续, 概率密度 $P(p)=|a_2\widetilde w_2(p/\hbar)+a_4\widetilde w_4(p/\hbar)|^2/\hbar$, 其中
\begin{equation}
 \widetilde w_n(k)=\sqrt{\frac L\pi}\e^{\ii kL/2}
 \frac{n\pi[1-(-1)^n\e^{-\ii kL}]}{(n\pi)^2-(kL)^2}.
\end{equation}
有限区间正弦的 Fourier 谱不是只有 $\pm n\pi/L$ 两根 delta. 后者只适用于全空间平面波叠加, 不能直接用于硬墙态.
\end{solution}
\begin{exercise}\label{ex:2202-4-4}
原题 4. 三维无限阱 $0<x<L_x$, $0<y<L_y$, $0<z<L_z$. (a) 解能量本征问题和边界条件. (b) 需要几个量子数?
\end{exercise}
\begin{solution}
分离变量给 $\psi=\sqrt{8/(L_xL_yL_z)}\prod_i\sin(n_i\pi x_i/L_i)$,
$E=\hbar^2\pi^2[n_x^2/L_x^2+n_y^2/L_y^2+n_z^2/L_z^2]/(2m)$, 三个独立正整数量子数. 立方盒的排列可有简并, 但态仍须三个标签, 不能只凭总能区分简并态.
\end{solution}
\begin{exercise}\label{ex:2202-4-5}
原题 5. 对称双阱: 两井分别 $-b-w<x<-b$ 与 $b<x<b+w$, 势 $-V_H$; 外边界无限. 图 A 的中央区势无限, 图 B 的中央 $|x|<b$ 势零. (a) 图 A 的定态及能级? (b) 图 B 各区波形, $b\to0,w\to0$ 极限? (c) 对称性条件? (d) 写图 B 的偶宇称负能解及本征值条件.
\end{exercise}
\begin{solution}
(a) 两个独立宽 w 阱, 每边 $\sin[n\pi(x-x_{\rm left})/w]$, 能量 $E_n=-V_H+\hbar^2\pi^2n^2/(2mw^2)$, 左右两重简并, 可取偶奇组合. (b) 对 $-V_H<E<0$, 阱内振荡, 中央为双曲函数, 外侧零. 中央有限势垒使左右耦合, 简并劈裂. $b\to0$ 成为宽 $2w$ 的单无限阱, 底势 $-V_H$; $w\to0$ 吸引井消失, 留宽 $2b$、底势零的无限阱, 不再有负能态.

(c) 取偶或奇宇称, 偶态在零处导数零, 奇态在零处函数零. (d) 记物理能量 $E=-\varepsilon$, $k^2=2m(V_H-\varepsilon)/\hbar^2$, $\kappa^2=2m\varepsilon/\hbar^2$. 偶态中央 $A\cosh\kappa x$, 右井 $D\sin[k(b+w-x)]$, 左井由偶延拓. 在 x=b 匹配得 $\kappa\tanh(\kappa b)=-k\cot(kw)$. 奇态则 $\kappa\coth(\kappa b)=-k\cot(kw)$. 这些条件加归一化决定实际波形, 不把随意正弦示意当精确解.
\begin{center}\begin{tikzpicture}[x=1cm,y=0.65cm]
\draw[->](-3,0)--(3,0)node[right]{$x$};\draw[thick](-2.5,1.4)--(-2.5,-1)--(-1,-1)--(-1,0)--(1,0)--(1,-1)--(2.5,-1)--(2.5,1.4);
\node[below]at(-1,0){$-b$};\node[below]at(1,0){$b$};\node[below]at(0,-1.4){中央有限势垒, 两侧井深 $V_H$};
\end{tikzpicture}\end{center}
\begin{center}\begin{tikzpicture}[x=1cm,y=.8cm]
\draw[->](-2.8,0)--(2.8,0)node[right]{$x$};
\draw[dashed](-2.5,0)--(-2.5,1.6);\draw[dashed](2.5,0)--(2.5,1.6);
\draw[dashed](-1,0)--(-1,1.6);\draw[dashed](1,0)--(1,1.6);
\draw[structurecolor,thick,domain=-1:1,samples=60]plot(\x,{.4*cosh(\x)});
\draw[structurecolor,thick,domain=1:2.5,samples=60]plot(\x,{.68993*sin(deg(1.352351*(2.5-\x)))});
\draw[structurecolor,thick,domain=-2.5:-1,samples=60]plot(\x,{.68993*sin(deg(1.352351*(2.5+\x)))});
\node[above]at(0,.8){偶态: 中央双曲余弦};\node[above]at(-1.8,1.4){井内振荡};\node[above]at(1.8,1.4){井内振荡};
\node[below]at(-2.5,0){$-b-w$};\node[below]at(-1,0){$-b$};\node[below]at(1,0){$b$};\node[below]at(2.5,0){$b+w$};
\end{tikzpicture}\end{center}
图为负能偶态的形态示意; 中央最小、两井振荡并在硬墙归零. 奇态中央改为 sinh, 两井符号相反; 真实能量仍由前述匹配条件确定.

\end{solution}
\begin{exercise}\label{ex:2202-4-6}
原题 6. 有限阱宽 $2a$, 题称``刚束缚的奇宇称基态'', 波长 $\lambda=4a$. (a) 算 $E_{\rm kin}=\hbar^2k^2/(2m)$ 和结合能, 图示; 给 $a=6$ fm、$V_H=5$ MeV、$mc^2=1876$ MeV. (b) $a=1.5$ fm 时如何?
\end{exercise}
\begin{solution}
按指定波长作估算, $k=\pi/(2a)$, $E_{\rm kin}=\hbar^2\pi^2/(8ma^2)$. 用 $\hbar c=197$ MeV fm, $a=6$ 时约 $0.709$ MeV, $E\simeq-4.291$ MeV, 结合能约 $4.291$ MeV. 阱底 $-5$、估计态能 $-4.291$、连续谱阈值 $0$ 即构成能量图. 宽度缩至四分之一, 估计动能增16倍为 $11.34$ MeV, 大于井深, 该指定波形不再给负能.

源表述有两处限制: 一维有限对称阱真基态是偶宇称, $\lambda=4a$ 也不是其精确边界解; 若指径向 u 的奇延拓, 需另给半直线条件与约化质量. 上述数值仅是题设波长估算, 不称作精确氘核束缚能, 也不由这一估算证明一维浅阱完全无束缚态. 一维任意窄的吸引方阱仍可有很浅的偶基态.

\begin{center}\begin{tikzpicture}[x=1cm,y=.8cm]
\draw[->](-2.7,0)--(3.6,0)node[right]{$x$};\draw[thick](-2.6,0)--(-1.5,0)--(-1.5,-2.5)--(1.5,-2.5)--(1.5,0)--(3.3,0);
\draw[structurecolor,dashed](-1.5,-2.1455)--(2.5,-2.1455)node[right]{$E\simeq-4.291$ MeV};
\draw[<->](-.6,-2.5)--(-.6,-2.1455);\node[below]at(0,-2.7){$E_{\rm kin}\simeq.709$ MeV};
\draw[<->](2,-2.1455)--(2,0)node[midway,right]{$B\simeq4.291$ MeV};
\node[above]at(0,0){连续谱阈值 $0$};\node[left]at(-1.5,-2.5){$-5$ MeV};
\end{tikzpicture}\end{center}
此能线按指定 $\lambda=4a$ 估计, 不表示有限阱精确本征解.
\end{solution}
